% Seminar 2, pagina 1 - manuscris datat 10-oct.-1994.
% Ultima cifră a anului este parţial acoperită de colţul deteriorat al foii;
% anul 1994 este dedus din cronologia notiţelor.
% Acesta este un fragment pentru \input, fără preambul şi fără \clearpage.
% \subsection*{Seminar 2 - Pagina 1}
\phantomsection
\subsection*{Seminar 2}
\addcontentsline{toc}{subsection}{Seminar 2}

\begin{center}
% \underline{\textbf{ALGEBRA}}\\[2pt]
% \textit{(Seminar 2)}
\end{center}
\begin{flushright}
\textit{10-oct.-1994}
\end{flushright}

\noindent Luni:\quad seminar alg. 3 ore;\quad seminar geom. 3 ore.

\bigskip
\noindent\fbox{1}\quad Fie $n\in\mathbb{N}^{*}$. Să se calculeze
\[
 d(n)=\left|\left\{
 f:\{1,2,\ldots,n\}\longrightarrow\{e_1,e_2,-e_1,-e_2\}
 \ \middle|\ \sum_{k=1}^{n}f(k)=0
 \right\}\right|,
\]
\noindent unde $e_1=(1,0)$, $e_2=(0,1)$.

\medskip
\noindent\underline{Soluţie:}
\[
 d(n)=\begin{cases}
 0, & n\text{ impar},\\[2pt]
 \text{de calculat}, & n\text{ par}.
 \end{cases}
\]

\noindent Fie $n=2m$. Atunci
\[
 d(2m)=\sum_{j=0}^{m}
 \binom{2m}{j}\binom{2m-j}{j}\binom{2m-2j}{m-j}.
\]
\[
\begin{aligned}
 d(2m)
 &=\sum_{j=0}^{m}
 \frac{(2m)!}{j!(2m-j)!}\,
 \frac{(2m-j)!}{j!(2m-2j)!}\,
 \frac{(2m-2j)!}{(m-j)!(m-j)!}\\[4pt]
 &=\sum_{j=0}^{m}\frac{(2m)!}{(m!)^2}
 \left(\frac{m!}{j!(m-j)!}\right)^2\\[4pt]
 &=\frac{(2m)!}{(m!)^2}\sum_{j=0}^{m}\binom{m}{j}^{2}\\[4pt]
 &=\frac{(2m)!}{(m!)^2}\binom{2m}{m}\\[4pt]
 &=\frac{(2m)!}{(m!)^2}\cdot\frac{(2m)!}{m!\,m!}
 =\frac{(n!)^2}{\left[\left(\frac n2\right)!\right]^4}.
\end{aligned}
\]
% În manuscris apare lângă prima sumă o schiţă de repartizare a celor 2m
% poziţii, cu etichetele j şi 2m-2j; formula combinatorică este păstrată.

\bigskip
\noindent\fbox{2}\quad\underline{(Caracterizare fct. inj.)}

\noindent Fie $f:A\longrightarrow B$. Să se arate echivalenţa:
\begin{enumerate}
\item[a)] $f$ e injectivă.
\item[b)] $f$ are o retractă, adică există $r:B\longrightarrow A$
 a.î.\ $r\circ f=1_A$.
\item[c)] $f$ este monomorfism: pentru orice mulţime $X$ şi
 $f_1,f_2:X\longrightarrow A$,
 \[
 f\circ f_1=f\circ f_2\ \Longrightarrow\ f_1=f_2.
 \]
 \[
 \begin{tikzcd}[column sep=large]
 X \arrow[r, shift left=.55ex, "f_1"]
   \arrow[r, shift right=.55ex, "f_2"']
 & A \arrow[r, "f"] & B
 \end{tikzcd}
 \]
\end{enumerate}
% Nota marginală uneşte grafic (a) şi (c); nu este un pas de demonstraţie.
% Enunţul (a) <=> (b) presupune A nevidă (ori o convenţie corespunzătoare
% pentru cazurile cu mulţimi vide); această condiţie nu apare în manuscris.
